%
% tab_bar_selection.tex
%
% A Z specification of the lux tab_bar active-tab selection state machine:
% the per-frame fire / echo / honour discipline that decides when a
% ``TabChanged'' event is emitted.  Formalises the design hardened
% empirically across four review rounds of PR lux-4n5n.
%
% Source of record:
%   src/punt_lux/display/renderers/imgui/tab_bar.py
%       -- ImGuiTabBarRenderer: begin_tab (the fire decision + force-select),
%          end (the honoured / pending writes)
%   src/punt_lux/scene/widget_state.py
%       -- WidgetState: reset_honoured (the re-push reset),
%          discard_for (the removal clear)
%   src/punt_lux/scene/manager.py
%       -- SceneManager._replace_scene_state: drives reset_honoured on re-push
%
% Type-check:  fuzz -t docs/tab_bar_selection.tex
% Model-check (see Section "Verification" for the exact goal predicates):
%   probcli docs/tab_bar_selection.tex -model_check -p DEFAULT_SETSIZE 2
%   probcli docs/tab_bar_selection.tex -model_check -nodead \
%       -goal "spurious=ztrue" -p DEFAULT_SETSIZE 2
%   probcli docs/tab_bar_selection.tex -model_check -nodead \
%       -goal "fires>1" -p DEFAULT_SETSIZE 2
%   probcli docs/tab_bar_selection.tex -model_check -nodead \
%       -goal "honoured/={} & selected/=active & clicked=zfalse" \
%       -p DEFAULT_SETSIZE 2
%
% Formatting follows the fuzz house rules: \quad~ for continuation
% indent (never \t1..\t9), schema lines under 80 columns, predicates
% broken at \land/\lor/\implies with the operator starting the
% continuation line.
%

\documentclass[a4paper,10pt,fleqn]{article}
\usepackage[margin=1in]{geometry}
\usepackage{fuzz}
\usepackage[colorlinks=true,linkcolor=blue,citecolor=blue,urlcolor=blue]{hyperref}

\begin{document}

\title{The lux Tab-Bar Selection Discipline: a Z Specification}
\author{Formal model of the fire / echo / honour state machine}
\date{July 2026}
\maketitle

% ============================================================
\section{Introduction}
% ============================================================

A \texttt{tab\_bar} carries one Hub-authoritative selection --- the
\emph{active} tab --- and is rendered every frame by an ImGui tab bar that
keeps its \emph{own} idea of which tab is selected.  The renderer must
reconcile the two: it must \emph{honour} the Hub's choice (force-select the
declared active tab so ImGui's default of tab~0 cannot clobber it), and it
must \emph{fire} a \texttt{TabChanged} event when, and only when, a genuine
user click moves ImGui's selection away from the active tab.  An echo --- the
Hub's own value coming back round on the next frame --- must fire nothing, or
a \texttt{fire $\to$ Hub $\to$ re-push $\to$ fire} loop would run.

The renderer arbitrates with a small piece of per-render-session bookkeeping
held in \texttt{WidgetState}: the \emph{honoured} value, the active tab the
frame last force-selected.  Before anything is honoured the slot is empty
(\texttt{\_UNHONOURED}); \texttt{end()} writes the current active tab into it
each frame; a whole-scene re-push resets it (\texttt{reset\_honoured}), and a
removal clears it (\texttt{discard\_for}).

The honoured value alone is not enough.  It records which \emph{Hub} value was
honoured; it does not record which \emph{user} selection has already been
reported.  In the window between a user's click and the Hub's re-push --- the
click-to-re-push latency window --- ImGui keeps reporting the clicked tab as
selected while the active tab has not yet moved.  With the honoured value
equal to the active tab, every frame in that window reads as a fresh user
switch and re-fires.  Closing that defect is what forces a \emph{second} slot,
the \emph{pending} selection: the tab for which a \texttt{TabChanged} is
already outstanding.  This document derives that slot and proves the four
safety properties the design must have.

This is a finite state machine over one tab bar and a small fixed set of tab
ids, so ProB can enumerate every interleaving of render frames, user clicks,
re-pushes, and removals, and either confirm the invariants or produce the
exact offending trace.

% ============================================================
\section{Basic Types}
% ============================================================

\begin{zed}
  [TABID]
\end{zed}

\texttt{TABID} is the set of tab identities the bar may hold.  Two ids suffice
to exhibit every defect: a declared active tab distinct from ImGui's tab-0
default, and a click that moves the selection between them.  Three exercise the
click-away-and-back paths more thoroughly.  The carrier is bounded by ProB's
\texttt{DEFAULT\_SETSIZE}, so no explicit cardinality constant is needed.

% ============================================================
\section{Free Types and Constants}
% ============================================================

\begin{zed}
  ZBOOL ::= ztrue | zfalse
\end{zed}

A two-valued flag written as a free type rather than the toolkit
\texttt{\bool}, so ProB animates it as a small enumerated set.

\begin{zed}
  MAXFIRE == 2
\end{zed}

The fire counter saturates at \texttt{MAXFIRE}.  A single click may fire at
most once; a count of two is already a witness to the repeated-fire defect, so
counting no higher keeps the state space small without hiding the violation.

\begin{axdef}
  home : TABID
\end{axdef}

\texttt{home} is ImGui's default selection, tab~0.  A freshly rendered tab bar
starts with ImGui showing \texttt{home}; the declared active tab may be some
\emph{other} id, and the first-frame honour must force the declared tab over
this default.

% ============================================================
\section{State}
% ============================================================

The state is flat: one schema, seven components.  Its predicate carries only
\emph{structural} invariants --- that the two bookkeeping slots hold at most
one id and the fire counter stays in range.  These hold in every design,
correct or buggy, so they act as coherence constraints, not as the safety
properties under test.  The safety properties (Section~\ref{sec:inv}) are
deliberately \emph{absent} here, so that a violating state remains reachable
and ProB can find it.

\begin{schema}{TabBar}
  active : TABID \\
  selected : TABID \\
  honoured : \power TABID \\
  pending : \power TABID \\
  clicked : ZBOOL \\
  fires : 0 \upto MAXFIRE \\
  spurious : ZBOOL
\where
  \# honoured \leq 1 \\
  \# pending \leq 1
\end{schema}

\texttt{active} is the Hub-authoritative selection.  \texttt{selected} is the
tab ImGui currently reports open.  \texttt{honoured} is the echo-suppression
slot: empty for \texttt{\_UNHONOURED}, or the singleton of the active tab the
last frame force-selected.  \texttt{pending} is the fire-suppression slot:
empty when no fire is outstanding, or the singleton of the tab a
\texttt{TabChanged} has already been fired for and not yet ratified by a
re-push.  Modelling the two slots as sets of size at most one lets the empty
set stand for ``none'' without a sentinel value.

The remaining three components are pure bookkeeping for the invariant checks,
not part of the renderer's own state.  \texttt{clicked} records whether the
current divergence between \texttt{selected} and \texttt{active} was produced
by a genuine user gesture; a fire raised while \texttt{clicked = zfalse} is
spurious.  \texttt{fires} counts fires since the last click or re-push --- the
window in which at most one fire is legitimate.  \texttt{spurious} latches
\texttt{ztrue} the moment any spurious fire occurs.

% ============================================================
\section{Initialization}
% ============================================================

\begin{schema}{Init}
  TabBar'
\where
  active' = home \\
  selected' = home \\
  honoured' = \emptyset \\
  pending' = \emptyset \\
  clicked' = zfalse \\
  fires' = 0 \\
  spurious' = zfalse
\end{schema}

A single initial state: nothing honoured, nothing pending, ImGui and the Hub
both on \texttt{home}, no fire seen.

% ============================================================
\section{Hub-Driven Operations}
% ============================================================

Three operations set the active tab from outside the renderer: the client
declaring the bar, a whole-scene re-push of a surviving bar, and a removal
followed by re-adding a same-id bar.  All three reset the echo- and
fire-suppression slots, because a fresh Hub value must be re-honoured from
scratch.  This reset is the guard that Section~\ref{sec:fidelity} removes to
reproduce three of the four defects.

\texttt{Declare} models the client installing the bar with a declared active
tab.  ImGui starts a fresh widget, so \texttt{selected} is its tab-0 default,
\texttt{home}; the declared active tab may differ.

\begin{schema}{Declare}
  \Delta TabBar \\
  v? : TABID
\where
  active' = v? \\
  selected' = home \\
  honoured' = \emptyset \\
  pending' = \emptyset \\
  clicked' = zfalse \\
  fires' = 0 \\
  spurious' = spurious
\end{schema}

\texttt{Repush} models a whole-scene re-push of a bar that \emph{survives} the
re-push (\texttt{SceneManager.\_replace\_scene\_state} calling
\texttt{reset\_honoured}).  The ImGui widget persists, so its \texttt{selected}
is kept; only the Hub value and the suppression slots change.

\begin{schema}{Repush}
  \Delta TabBar \\
  v? : TABID
\where
  active' = v? \\
  selected' = selected \\
  honoured' = \emptyset \\
  pending' = \emptyset \\
  clicked' = zfalse \\
  fires' = 0 \\
  spurious' = spurious
\end{schema}

\texttt{RemoveReadd} models the bar being removed and a same-id bar re-added
(\texttt{discard\_for} clearing the slots).  Its state effect is identical to
\texttt{Repush} --- a stale ImGui selection kept, the slots cleared --- but it
travels a different code path (\texttt{discard\_for}, not
\texttt{reset\_honoured}).  The two operations are kept distinct because each
guards invariant~1 on a path the other does not cover; removing either guard
alone must still be caught.

\begin{schema}{RemoveReadd}
  \Delta TabBar \\
  v? : TABID
\where
  active' = v? \\
  selected' = selected \\
  honoured' = \emptyset \\
  pending' = \emptyset \\
  clicked' = zfalse \\
  fires' = 0 \\
  spurious' = spurious
\end{schema}

% ============================================================
\section{The User Click}
% ============================================================

A user clicks a tab other than the active one.  ImGui moves its selection;
the click opens the latency window (\texttt{clicked} becomes \texttt{ztrue})
and starts a fresh fire count.  The click itself fires nothing --- firing is
the render frame's decision.

\begin{schema}{UserClick}
  \Delta TabBar \\
  u? : TABID
\where
  u? \neq active \\
  selected' = u? \\
  clicked' = ztrue \\
  fires' = 0 \\
  active' = active \\
  honoured' = honoured \\
  pending' = pending \\
  spurious' = spurious
\end{schema}

% ============================================================
\section{The Render Frame}
% ============================================================

One frame runs \texttt{begin\_tab} over every tab and then \texttt{end()}.
Exactly one of three cases applies, so a render step is \emph{always} enabled
--- the reason the machine never deadlocks (Section~\ref{sec:inv},
invariant~4).  Every case ends by writing the active tab into
\texttt{honoured}: that is \texttt{end()}'s honoured write, and it is what
makes the \emph{next} frame read the current active tab as ``already
honoured''.

\texttt{FrameForceSelect} is the fresh-value case: the active tab is not the
honoured one, so \texttt{begin\_tab} force-selects it.  ImGui's selection is
driven to the active tab, and no fire can occur --- the only tab reported
selected is the active tab itself, which is never a switch.  This is both the
first-frame honour and the echo suppression after a Hub-driven change.

\begin{schema}{FrameForceSelect}
  \Delta TabBar
\where
  active \notin honoured \\
  selected' = active \\
  honoured' = \{ active \} \\
  pending' = pending \\
  clicked' = zfalse \\
  fires' = fires \\
  spurious' = spurious \\
  active' = active
\end{schema}

\texttt{FrameFire} is the genuine-switch case: the active tab is already
honoured (no force-select this frame), ImGui reports a \emph{different} tab
selected, and no fire is yet outstanding for that tab.  The frame fires,
records the tab in \texttt{pending}, and bumps the counter.  A fire raised
while \texttt{clicked = zfalse} --- no user gesture behind it --- latches
\texttt{spurious}.

\begin{schema}{FrameFire}
  \Delta TabBar
\where
  active \in honoured \\
  selected \neq active \\
  selected \notin pending \\
  selected' = selected \\
  honoured' = \{ active \} \\
  pending' = \{ selected \} \\
  clicked' = clicked \\
  fires' = \IF fires < MAXFIRE \THEN fires + 1 \ELSE MAXFIRE \\
  spurious' = \IF clicked = zfalse \THEN ztrue \ELSE spurious \\
  active' = active
\end{schema}

The guard \texttt{selected $\notin$ pending} is the fix.  Once a tab is
pending, later frames in the same window see it in \texttt{pending} and fall
to \texttt{FrameSteady} instead of re-firing.  Deleting that conjunct (and the
\texttt{pending' = \{selected\}} write that feeds it) is the buggy variant of
Section~\ref{sec:fidelity}, defect~(d).

\texttt{FrameSteady} is every other honoured frame: ImGui is on the active tab,
or on a tab whose fire is already outstanding.  Nothing fires; the honoured
write still happens.

\begin{schema}{FrameSteady}
  \Delta TabBar
\where
  active \in honoured \\
  (selected = active \lor selected \in pending) \\
  selected' = selected \\
  honoured' = \{ active \} \\
  pending' = pending \\
  clicked' = clicked \\
  fires' = fires \\
  spurious' = spurious \\
  active' = active
\end{schema}

The three frame guards partition the state space: \texttt{active $\notin$
honoured} selects \texttt{FrameForceSelect}; \texttt{active $\in$ honoured}
splits on the fire condition into \texttt{FrameFire} and \texttt{FrameSteady}.
Exactly one is enabled in every state, so rendering never stalls.

% ============================================================
\section{Invariants}
\label{sec:inv}
% ============================================================

Four safety properties must hold across all reachable states, plus
deadlock-freedom.  They are stated here as predicates over \texttt{TabBar};
how each is discharged is given in Section~\ref{sec:verify}.  None is placed in
the \texttt{TabBar} schema: a predicate placed there would be enforced as a
guard on every successor, silently disabling any operation that would break it
and thereby masking the very defect this model exists to catch.

\paragraph{Invariant 1 --- no spurious fire.}
A \texttt{TabChanged} is fired only on a genuine user switch, never off a Hub
echo or a stale honoured value.  Equivalently, \texttt{spurious} is never set:
\[
  spurious = zfalse.
\]
This subsumes echo suppression --- a Hub-driven change (\texttt{Declare},
\texttt{Repush}, \texttt{RemoveReadd}) carries \texttt{clicked = zfalse}, so
any fire it provokes is spurious and would trip this invariant.

\paragraph{Invariant 2 --- no repeated fire.}
A single user click yields at most one fire across all frames until the
re-push.  With the counter reset on every click and re-push:
\[
  fires \leq 1.
\]

\paragraph{Invariant 3 --- first-frame honour.}
Once a frame has rendered (\texttt{honoured} is non-empty), if the user has not
clicked then ImGui shows the Hub's active tab --- the declared active tab is
never left clobbered by ImGui's tab-0 default:
\[
  honoured \neq \emptyset \land clicked = zfalse \implies selected = active.
\]

\paragraph{Invariant 4 --- deadlock freedom.}
No reachable state leaves the machine unable to progress.  The three frame
cases partition the state space, so a render step is enabled in every state.

\paragraph{Echo suppression.}
A Hub-driven active-tab change fires nothing.  This is not a separate
reachability goal: \texttt{Declare}, \texttt{Repush}, and \texttt{RemoveReadd}
never fire, and each resets \texttt{honoured} to empty, so the frame that
follows is always \texttt{FrameForceSelect}, which cannot fire.  A failure to
suppress an echo is therefore a spurious fire, caught by invariant~1.  The
fidelity traces for defects (b) and (c) are exactly echo-suppression failures.

% ============================================================
\section{Verification}
\label{sec:verify}
% ============================================================

Type-checking is a \texttt{fuzz} pass:
\begin{verbatim}
  fuzz -t docs/tab_bar_selection.tex
\end{verbatim}

Invariants 1, 2, and 3 are state properties, checked by asking ProB whether
their negation is \emph{reachable}.  A goal that is not found means the
invariant holds on every reachable state:
\begin{verbatim}
  # Invariant 1 (must report: goal NOT found)
  probcli docs/tab_bar_selection.tex -model_check -nodead \
    -goal "spurious=ztrue" -p DEFAULT_SETSIZE 2

  # Invariant 2 (must report: goal NOT found)
  probcli docs/tab_bar_selection.tex -model_check -nodead \
    -goal "fires>1" -p DEFAULT_SETSIZE 2

  # Invariant 3 (must report: goal NOT found)
  probcli docs/tab_bar_selection.tex -model_check -nodead \
    -goal "honoured/={} & selected/=active & clicked=zfalse" \
    -p DEFAULT_SETSIZE 2
\end{verbatim}

Invariant 4 is the deadlock check over the full reachable state space:
\begin{verbatim}
  # Invariant 4 (must report: no deadlock, no invariant violation)
  probcli docs/tab_bar_selection.tex -model_check -p DEFAULT_SETSIZE 2
\end{verbatim}

% ============================================================
\section{Fidelity: reproducing the four defects}
\label{sec:fidelity}
% ============================================================

A model that cannot reproduce the known defects proves nothing.  Each defect is
the current specification with one guard removed; each then makes the
corresponding invariant's negation reachable.  Below, each variant names the
edit, the goal that must change from \emph{not found} to \emph{found}, and the
shortest offending trace.  Under \texttt{DEFAULT\_SETSIZE 2},
\texttt{TABID = \{TABID1, TABID2\}} with \texttt{home = TABID1}; write
$t_1 = home$ and $t_2$ for the other id.

\paragraph{Defect (a) --- first-frame honoured defaulting to active.}
Edit: \texttt{Declare} seeds \texttt{honoured' = \{v?\}} instead of
\texttt{$\emptyset$}, modelling the original code reading the honoured slot's
default as \texttt{active} rather than \texttt{\_UNHONOURED}.  Goals found:
both invariant~3 (\texttt{honoured/=\{\} \& selected/=active \&
clicked=zfalse}) and invariant~1 (\texttt{spurious=ztrue}).  Trace:
\begin{enumerate}
  \item \texttt{Init}: \texttt{active = selected = $t_1$}.
  \item \texttt{Declare($t_2$)} in the buggy form: \texttt{active = $t_2$},
    \texttt{selected = $t_1$}, \texttt{honoured = \{$t_2$\}},
    \texttt{clicked = zfalse}.  This state alone violates invariant~3:
    \texttt{honoured} is non-empty, the user has not clicked, yet
    \texttt{selected = $t_1$ $\neq$ $t_2$ = active} --- ImGui's tab-0 default
    stands where the declared tab should be.
  \item \texttt{FrameFire}: \texttt{active = $t_2$ $\in$ honoured},
    \texttt{selected = $t_1$ $\neq$ $t_2$}, \texttt{$t_1$ $\notin$ pending}, so
    the frame fires with \texttt{clicked = zfalse} --- \texttt{spurious} is set,
    violating invariant~1.
\end{enumerate}
Under the intact guard, \texttt{Declare($t_2$)} leaves \texttt{honoured =
$\emptyset$}, so the first frame is \texttt{FrameForceSelect}: it drives
\texttt{selected} to $t_2$ and fires nothing.

\paragraph{Defect (b) --- honoured surviving a re-push.}
Edit: \texttt{Repush} omits \texttt{honoured' = $\emptyset$}, keeping
\texttt{honoured' = honoured} (the original \texttt{reset\_honoured} not
running on the re-push path).  Goal found: invariant~1
(\texttt{spurious=ztrue}).  Trace:
\begin{enumerate}
  \item \texttt{Init}; \texttt{FrameForceSelect}: \texttt{honoured = \{$t_1$\}},
    \texttt{selected = active = $t_1$}.
  \item \texttt{UserClick($t_2$)}: \texttt{selected = $t_2$},
    \texttt{clicked = ztrue}.
  \item \texttt{FrameFire}: fires once (legitimately), \texttt{pending =
    \{$t_2$\}}, \texttt{honoured = \{$t_1$\}}.
  \item \texttt{Repush($t_1$)} in the buggy form --- a re-push that does not
    change this bar's active tab: \texttt{active = $t_1$}, \texttt{honoured}
    kept as \texttt{\{$t_1$\}}, \texttt{pending = $\emptyset$},
    \texttt{clicked = zfalse}, \texttt{selected = $t_2$} kept.
  \item \texttt{FrameFire}: \texttt{active = $t_1$ $\in$ honoured},
    \texttt{selected = $t_2$ $\neq$ $t_1$}, \texttt{$t_2$ $\notin$ pending}, so
    it fires with \texttt{clicked = zfalse} --- \texttt{spurious} is set.  The
    stale honoured value made the re-push's own selection read as a fresh user
    switch.
\end{enumerate}
Under the intact guard, \texttt{Repush} empties \texttt{honoured}, so the next
frame is \texttt{FrameForceSelect}: it drives \texttt{selected} back to the
Hub's $t_1$ and fires nothing.

\paragraph{Defect (c) --- honoured not cleared on removal.}
Edit: \texttt{RemoveReadd} omits \texttt{honoured' = $\emptyset$} (the original
\texttt{discard\_for} not clearing the honoured key).  Goal found:
invariant~1.  The trace is that of defect~(b) with \texttt{RemoveReadd($t_1$)}
in place of \texttt{Repush($t_1$)} at step~4: a same-id bar re-added while its
stale honoured value survives reads its stale ImGui selection as a fresh switch
and fires spuriously on re-add.  Under the intact guard, the re-add empties
\texttt{honoured} and the next frame force-selects without firing.

\paragraph{Defect (d) --- honoured not advanced on a user switch (the OPEN
finding).}
Edit: \texttt{FrameFire} drops the \texttt{selected $\notin$ pending} guard and
the \texttt{pending' = \{selected\}} write, modelling the original renderer
that wrote only the active tab into the honoured slot and kept no record of
the tab it had already fired for.  Goal found: invariant~2
(\texttt{fires>1}).  Trace:
\begin{enumerate}
  \item \texttt{Init}; \texttt{FrameForceSelect}: \texttt{honoured = \{$t_1$\}},
    \texttt{selected = active = $t_1$}.
  \item \texttt{UserClick($t_2$)}: \texttt{selected = $t_2$},
    \texttt{clicked = ztrue}, \texttt{fires = 0}.
  \item \texttt{FrameFire}: \texttt{active = $t_1$ $\in$ honoured},
    \texttt{selected = $t_2$ $\neq$ $t_1$}, so it fires --- \texttt{fires = 1}.
    In the buggy form nothing records $t_2$ as pending.
  \item \texttt{FrameFire} again --- the next frame in the same latency window,
    before any re-push, with \texttt{active = $t_1$}, \texttt{honoured =
    \{$t_1$\}}, \texttt{selected = $t_2$} unchanged: the guard that would have
    seen $t_2$ pending is gone, so it fires \emph{again} --- \texttt{fires = 2},
    violating invariant~2.  Every further frame in the window re-fires.
\end{enumerate}
Under the intact guard, step~3 records \texttt{pending = \{$t_2$\}}, so step~4
is \texttt{FrameSteady}: \texttt{selected = $t_2$ $\in$ pending}, no fire.  The
window fires exactly once.

\paragraph{Summary.}
Each defect's guard, removed, makes its invariant's negation reachable; with
every guard intact, all four reachability goals return \emph{not found} and the
deadlock check reports none.  That difference is the evidence that the model
detects the defects the code was changed to prevent --- and, for defect~(d),
that the \texttt{pending} slot is the correct fix: it is the minimal state that
makes the latency window fire exactly once.

% ============================================================
\section{From the Model to the Code}
% ============================================================

The model's \texttt{pending} slot is a second echo-suppression key in
\texttt{WidgetState}, beside the honoured key.  \texttt{begin\_tab} reads it,
fires only when the reported selection is neither the active tab nor the
pending tab, and records the fired tab into it; \texttt{end()} is unchanged ---
it still writes only the active tab into the honoured slot.  The re-push and
removal resets (\texttt{reset\_honoured}, \texttt{discard\_for}) clear the
pending key alongside the honoured key, which is the model's \texttt{pending' =
$\emptyset$} in \texttt{Repush} and \texttt{RemoveReadd}.  No other honoured-
or fire-related behaviour changes.

\end{document}
